% page 580 of the facsimile (image p-579 of the scan)
% block 147.3 x 263.6 mm ; left margin 8.0 mm
\pgc{-12.700mm}{0.000mm}{
\l{}
\l{}
\l{}
\l{\cel{52}572}
\l{}
\l{+\sou{taxio}.\cel{2}Fiakro automobila, taximetroza. - DEFIRS.}
\l{}
\l{\sou{taxuso}. (\sou{bot.}) Arboro verda, ek la familio \textquotedbl{}koniferi\textquotedbl{}.}
\l{\cel{3}- L.\cel{1}\sou{taxus}. - DL}
\l{}
\l{\sou{tayo}. (\sou{anat.}) Parto streta di la homo-korpo, inter han-}
\l{\cel{3}chi e la infra parto di la busto. - FIS.}
\l{}
\l{\sou{teo}.\cel{2}(\sou{bot.}) I. Folii di arboreto (L.\cel{1}\sou{thea}) de Chinia,}
\l{\cel{3}de Japonia, qua, per infuzo, donas drinkajo aromat-}
\l{\cel{3}oza, stimuliva. - II. Ta ipsa infuzuro. - DEFIS.}
\l{}
\l{\sou{teatro}. I. Edifico ube reprezentesas verki qui havas}
\l{\cel{3}la formo di konversi. - II. Ta edifico kun ti qui}
\l{\cel{3}direktas lu, la aktori qui reprezentas la verki}
\l{\cel{3}teatrala, la dekoruri, e c. - DEFIRS.}
\l{}
\l{\sou{tebaido}.\cel{2}Retreteyo, refujeyo, solitara. - DFIS.}
\l{}
\l{\sou{tedar}. (\sou{trans.})\cel{3}Fatigar per arivir desoportune, o pro}
\l{\cel{3}tro repetar, tro iterar sua dici. - EIS.}
\l{}
\l{\sou{tegar}. (\sou{trans.}) Envelopar plu o min multa ek la parti (di}
\l{\cel{3}objekto). - L.}
\l{}
\l{\sou{tegulo}. Plaketo de terakoto, uzata por facar la tekti.}
\l{\cel{3}- IL.}
\l{}
\l{\sou{tegumento}. (anat.)\cel{2}Tisuo organika adheranta, qua tegas}
\l{\cel{3}ula parti di la animali, di la vejetanti. - DEFIS.}
\l{}
\l{\sou{teismo}.\cel{2}Kredo pri la existo di deo. - DEFIRS.}
\l{}
\l{\sou{teisto}. Persono qua kredas ke existas deo. - DEFIRS.}
\l{}
\l{\sou{teko}.\cel{2}(\sou{bot.)}\cel{2}Arboro di India, ek la familio \textquotedbl{}verbena-}
\l{\cel{3}cei\textquotedbl{}, di qua la ligno, extreme-densa, e harda, uzesas}
\l{\cel{3}por konstruktar domi, navi, e c. - L.\cel{1}\sou{tectonia}. DEFIRS}
\l{}
\l{\sou{tekniko.}\cel{2}Ensemblo ek la procedi di arto. - DEFIRS.}
\l{}
\l{\sou{tekenllogio}. Expliko di la procedi, di la termini qui kon-}
\l{\cel{3}cernas propre ta o ca arto o cienco. - DEFIRS.}
\l{}
\l{\sou{tekto}. (\sou{arkitekt.})\cel{2}Parto di la karpenturo somite di edi-}
\l{\cel{3}fico, di domo, ek teguli, ardezo-plaki, palio, e c.,}
\l{\cel{3}qua kovras ta konstrukturo. - FIS.}
\l{}
\l{\sou{telo}. Texuro de fili ek lino o kanabo.}
\l{}
\l{\sou{telefonar}. (\sou{trans., ad)}. Transmisar voce (ulo ad ulu) per}
\l{\cel{3}uzar instrumento ek plako dina de fero pura qua, kande}
\l{\cel{3}la voco-soni produktesas, vibras e vibrigas unisone,}
\l{\cel{3}per filo elektrala, lamo simila qua produktas, sur la}
\l{\cel{3}membrano di la timpano, vibri korespondanta e qui geni-}
\l{\cel{3}tas soni identa. - DEFIRS.}
}
